% concatenated from Convex_in_SCV_corrected_v2.tex
\documentclass[12pt]{amsart}

% --- Paquetes de idioma y codificación ---
\usepackage[utf8]{inputenc}      % Codificación UTF-8
\usepackage[T1]{fontenc}
\usepackage[english]{babel}
\usepackage{csquotes}

\usepackage{setspace}
% --- Paquetes de matemáticas ---
\usepackage{amsmath, amssymb, amsthm, mathtools}
\usepackage{mathrsfs}            % Letras caligráficas
\usepackage{esint}               % Integrales dobles, etc.

% --- Paquetes útiles ---
\usepackage{enumitem}            % Mejoras en listas enumeradas
\usepackage{graphicx}            % Imágenes
\usepackage{xcolor}              % Colores
\usepackage{hyperref}            % Hipervínculos
\hypersetup{
  colorlinks = true,
  linkcolor = blue,
  citecolor = blue,
  urlcolor  = blue
}

% --- Márgenes ---
\usepackage[margin=1.2in]{geometry}

% --- Teoremas y entornos matemáticos ---
\newtheorem{theorem}{Theorem}[section]
\newtheorem{lemma}[theorem]{Lemma}
\newtheorem{corollary}[theorem]{Corollary}
\newtheorem{proposition}[theorem]{Proposition}

\theoremstyle{definition}
\newtheorem{definition}[theorem]{Definition}

\theoremstyle{remark}
\newtheorem{remark}[theorem]{Remark}

\def\dis{\displaystyle}
\def\bq{\begin{equation}}
\def\eq{\end{equation}}
\def\bp{\begin{proof}}
\def\ep{\end{proof}}
\def\noi{\noindent}
\def\bi{\bigskip}
\def\me{\medskip}
\def\sm{\smallskip}
\def\P{{\mathbb P}^n}
\def\R{{\mathbb R}}
\def\S{{ S^{k}_{ij}}}
\def\SS{{\mathcal{S}}}
\def\D{{\mathbb D}}
\def\al{\alpha}
\def\va{\varphi}
\def\ga{\gamma}
\def\si{\sigma}
\def\om{\omega}
\def\C{{\mathbb C}}
\def\Cn{{\mathbb C}^n}
\def\la{{\lambda}}
\def\laa{{\lambda}_{\alpha}}
\def\mar{{\mu_r(\alpha)}}
\def\Fa{{\mathcal F}_{\alpha}}
\def\Far{{\mathcal F}_{\alpha, r}}
\def\Fao{{\mathcal F}_{\alpha}^0}
\def\I{\textsl{I}}
\newcommand{\norm}[1]{\left\Vert#1\right\Vert}
\newcommand{\abs}[1]{\left\vert#1\right\vert}
\newcommand{\set}[1]{\left\{#1\right\}}
\newcommand{\Real}{\mathbb R}
\newcommand{\eps}{\varepsilon}
\newcommand{\To}{\longrightarrow}
\newcommand{\BX}{\mathbf{B}(X)}
\newcommand{\A}{\mathcal{A}}

\def\P{{\mathbb P}}
\def\C{{\mathbb C}}
\def\B{{\mathbb B}}
\def\D{{\mathbb D}}
\def\R{{\mathbb R}}

% --- Numeración de ecuaciones ---
\numberwithin{equation}{section}

% --- Título y autor ---
\title[Schwarzian Derivative and Convexity]{The Schwarzian Derivative for Convex Holomorphic Mappings in Several Complex Variables}
\author{Rodrigo Hern\'andez}
\date{}

% --- Documento ---
\begin{document}

\begin{abstract}
We obtain upper bounds for the norm of the Schwarzian derivative of convex
holomorphic mappings defined on the polydisk and the unit ball in
$\mathbb{C}^n$. For coordinate-wise convex mappings on the polydisk, we derive
a sharp estimate extending the classical one-variable result of
Chuaqui--Duren--Osgood to higher dimensions. For the Roper--Suffridge
extension operator in the unit ball, we obtain an explicit bound that
represents the best available estimate in this setting. 
\end{abstract}

\subjclass[2020]{Primary 30C45, 32H02; Secondary 30C55, 30C80}
 
\keywords{Schwarzian derivative, convex mappings, polydisk, unit ball,
Roper--Suffridge operator, Bergman metric, several complex variables}
\maketitle


\section*{Introduction}

Few objects in complex analysis are as deceptively simple and as geometrically
rich as the Schwarzian derivative. Originally formalized by Nehari as a tool
for detecting univalence, it has since become a lens through which one can
read the geometry of holomorphic mappings. For a function $f$ analytic and
locally univalent in a domain of the complex plane, it is defined by
\[
S_f(z) = \left( \frac{f''(z)}{f'(z)} \right)' - \frac{1}{2}
\left( \frac{f''(z)}{f'(z)} \right)^2,
\]
and measures, in a precise sense, how far $f$ is from being a M\"{o}bius
transformation --- these are exactly the maps for which $S_f \equiv 0$.
Nehari's fundamental result \cite{nehari1949} states that the condition
\[
|S_f(z)| \leq \frac{2}{(1 - |z|^2)^2}
\]
is sufficient for the global univalence of $f$ in the unit disk $\mathbb{D}$,
with the constant $2$ being sharp.

A particularly striking instance of this theory arises for convex mappings.
Chuaqui, Duren, and Osgood showed that any holomorphic function mapping the
disk conformally onto a convex domain satisfies $\|S_f\| \leq 2$, and that
this bound is sharp \cite{chuaqui2011}. Their argument rests on a structural
representation of $S_f$ in terms of the Carath\'{e}odory family: there exists
$\varphi:\mathbb{D}\to\mathbb{D}$, with $\varphi(0)=0$, such that
\[
S_f(z) = \frac{2\varphi'(z)}{(1 - z\varphi(z))^2}.
\]
From this one deduces that
\[
\|S_f\| = \sup_{z \in \mathbb{D}} (1 - |z|^2)^2 |S_f(z)| \leq 2,
\]
with equality precisely when $f$ maps $\mathbb{D}$ onto a parallel strip.
Moreover, bounded convex mappings satisfy $\|S_f\| < 2$ strictly, since a
convex mapping attaining the value $2$ cannot map onto a quasidisk.

These ideas have had lasting influence on the study of function classes such
as convex functions of order $\alpha$, strongly convex functions, and
polygonal mappings, where sharp Schwarzian bounds encode quantitative
information about univalence and distortion \cite{kanas2011}. Yet, despite
this progress in one variable, the analogous question in higher dimensions ---
what are the sharp Schwarzian norm bounds for convex mappings in
$\mathbb{C}^n$? --- had remained unanswered.

Carrying this theory beyond one variable is not a straightforward task. In
$\mathbb{C}^n$ for $n > 1$, there is no canonical notion of conformality, and
the group of M\"{o}bius transformations no longer acts transitively on the
space of locally biholomorphic mappings. The Schwarzian derivative cannot
simply be a scalar invariant; its definition requires a more structural
reformulation, and several non-equivalent generalizations have been proposed
\cite{MM1, MM2, MT}, each recovering different facets of the classical theory.

In our previous work \cite{He, hernandez2007}, we introduced a Schwarzian
operator on domains in $\mathbb{C}^n$, formulated as a tensor of type $(1,2)$
that is invariant under linear changes of coordinates, and showed that norm
bounds on this operator imply univalence --- a higher-dimensional analogue of
Nehari's criterion. In a related direction, joint work with
Mart\'{\i}n \cite{hernandezmartin2013} extended the theory to complex harmonic
mappings, where we introduced harmonic analogues of the Schwarzian and
pre-Schwarzian, recovering results such as Becker's criterion in the
non-analytic setting.

The present paper takes up the convexity problem in this higher-dimensional
framework. Our goal is to establish explicit norm bounds for the Schwarzian
operator associated with convex mappings, in the spirit of
Chuaqui--Duren--Osgood. We work in two settings. For coordinate-wise convex
mappings on the polydisk $\mathbb{P}^n$, we obtain the sharp bound
\[
\|S_f\| \leq \frac{\sqrt{8(n+3)(n-1)}}{n+1}.
\]
For the unit ball $\mathbb{B}^n$, the situation is more delicate: convex
mappings are considerably harder to construct and parametrize than in the
polydisk, and the existing convexity criterion does not lend itself to explicit
norm calculations. The Roper--Suffridge operator offers a natural and elegant
path forward --- starting from a convex function in the disk, it produces a
convex mapping in the ball in a geometrically transparent way --- and within
this class we obtain the bound
\[
\|S_f\| \leq \frac{2}{3\sqrt{3(n+1)}}.
\]
Both estimates are consistent with the classical one-variable bound
$\|S_f\| \leq 2$ of Chuaqui--Duren--Osgood, to which they reduce in the
appropriate limiting sense.

\section{Preliminaries}\label{sec:prelim}

Let $f:\Omega\subset\Cn\rightarrow\Cn$ be a locally
biholomorphic mapping defined on some domain $\Omega$. Several non-equivalent
definitions of the Schwarzian derivative in several complex variables have
been proposed in the literature \cite{MM1, MM2, MT}; we work with the one
introduced in \cite{Oda}, which has the advantage of being tied directly to
the theory of completely integrable systems. Following \cite{Oda}, the
Schwarzian derivatives of $f=(f_1,\ldots,f_n)$ are defined as
\[
\S f= \dis\sum^{n}_{l=1}\frac{\partial^{2} f_{l}}{\partial z_{i}\partial
z_{j}} \frac{\partial z_{k}}{\partial
f_{l}}-\frac{1}{n+1}\left(\delta^{k}_{i}\frac{\partial}{\partial
z_{j}}+\delta^{k}_{j}\frac{\partial}{\partial z_{i}}\right)
\log(J_f)\, ,
\]
where $i,j,k=1,2, \ldots,n$ and $\delta^{k}_{i}$ are the Kronecker symbols.
For $n>1$, the Schwarzian derivatives satisfy $\S f=0$ for all
$i,j,k=1,2,\ldots,n$ if and only if $f(z)=M(z)$ for some M\"{o}bius
transformation
\[
M(z)=\left(\frac{l_1(z)}{l_0(z)},\ldots,\frac{l_n(z)}{l_0(z)}\right),
\]
where $l_i(z)=a_{i0}+a_{i1}z_1+\cdots+a_{in}z_n$ with $\det(a_{ij})\neq 0$.
Furthermore, for a composition one has
\bq\label{chain rule}
\S(g\circ f)(z)=\S f(z)+
\sum^{n}_{l,m,r=1}S^{r}_{lm}g(w)\frac{\partial w_{l}}{\partial z_{i}}
\frac{\partial w_{m}}{\partial z_{j}}
\frac{\partial z_{k}}{\partial w_{r}}\; ,\quad w=f(z)\, .
\eq
In particular, if $g$ is a M\"{o}bius transformation, then
$\S(g\circ f)=\S f$. The coefficients $S^0_{ij}f$ are given by
\[
S^0_{ij}f(z)=(J_f)^{\frac{1}{n+1}}\left(\frac{\partial^2}{\partial
z_i\partial z_j}(J_f)^{-\frac{1}{n+1}}-
\sum_{k=1}^n\frac{\partial}{\partial
z_k}(J_f)^{-\frac{1}{n+1}}S^k_{ij}f(z)\right).
\]
In \cite{Oda} one also finds a description of the functions with prescribed
Schwarzian derivatives $\S f$. Consider the following over-determined system
of partial differential equations:
\bq\label{system}
\frac{\partial^{2}u}{\partial z_{i}\partial z_{j}} =
\sum^{n}_{k=1}P^{k}_{ij}(z)\frac{\partial u}{\partial z_{k}}+
P^{0}_{ij}(z)u \; ,\quad i,j= 1,2, \ldots,n\, ,
\eq
where $z=(z_{1},z_{2},\ldots,z_n)\in \Omega$ and $P^k_{ij}(z)$ are holomorphic
functions in $\Omega$ for $i,j,k=0,\ldots,n$. The system \eqref{system} is
called \textit{completely integrable} if it admits $n+1$ (the maximum number
of) linearly independent solutions, and is said to be in \textit{canonical
form} (see \cite{yoshida1976}) if the coefficients satisfy
\[
\sum_{j=1}^{n} P_{ij}^{j}(z)=0\, ,\quad i=1,2,\ldots,n\, .
\]
It was shown that \eqref{system} is a completely integrable system in canonical
form if and only if $P^k_{ij}=\S f$ for a locally biholomorphic mapping
$f=(f_1, \ldots, f_n)$, where $f_i=u_i/u_0$ for $1\leq i\leq n$ and
$u_0, u_1,\ldots, u_n$ is a set of linearly independent solutions of the
system. For a given mapping $f$, the function $u_0=(J_f)^{-\frac{1}{n+1}}$ is
always a solution of \eqref{system} with $P_{ij}^k=\S f$.

We recall the following definitions from \cite{He}, in which the individual
Schwarzian derivatives $S^kf$ are combined to define the operators below.

\begin{definition}
For each $k=1,\ldots,n$, we let $S^{k}_f$ denote the $n\times n$ matrix
\[
S^k_f= (\S f)\, ,\quad i,j=1,\ldots,n\, .
\]
\end{definition}

\begin{definition}
Let $f:\Omega\rightarrow\Cn$ be locally biholomorphic, and let $T_z\Omega$
denote the tangent space at $z\in\Omega$. We define the \textit{Schwarzian
derivative operator} as the mapping
$S_f(z):T_z\Omega \to T_z\Omega$ given by
\[
S_f(z)(\vec{v},\vec{v})=\left(S^1_f(z)(\vec{v})\,,\ldots,S^{n}_f(z)\vec{v}\right),
\]
where $\vec{v}\in T_z \Omega$ and $S^k_f(z)(\vec{v})=\vec{v}^{\,t}S^k_f(z)\vec{v}$.
\end{definition}

\noi With this notation, equation \eqref{system} with $P^k_{ij}=\S f$ can be
rewritten as
\bq\label{Hess}
\mathrm{Hess}(u)(z)(\vec{v},\vec{v})
= S_f(z)(\vec{v})\cdot \nabla u(z)+S^0_f(z)(\vec{v})\,u\, ,
\eq
where $S^0_f(z)(\vec{v})=\vec{v}^{\,t}S^0_f(z)\vec{v}$.

In \cite{He} it is shown that
\begin{equation}\label{sf operator}
S_f(z)(\vec{v},\vec{v})=(Df(z))^{-1}D^2f(z)(\vec{v},\vec{v})
-\frac{2}{n+1}\bigl(\nabla\log J_f(z)\cdot\vec{v}\bigr)\,\vec{v}.
\end{equation} 

The Bergman metric in the polydisk
$\mathbb{P}^n=\{(z_1,\ldots,z_n): |z_i|<1,\,i=1,\ldots, n\}$
is the Hermitian metric defined by the diagonal matrix
\[
g_{ii}(z)=\frac{2}{(1-|z_i|^2)^2}\,.
\]
Let $\mathbb{B}^n=\{(z_1,\ldots,z_n)\in\mathbb{C}^n:|z_1|^2+\cdots+|z_n|^2<1\}$.
The Bergman metric on $\mathbb{B}^n$ is the Hermitian inner product defined by
\begin{equation}\label{bergman metric}
g_{ij}(z)=\frac{n+1}{(1-|z|^2)^2}\left[(1-|z|^2)\delta_{ij}+\bar{z}_iz_j\right],
\end{equation}
where $\delta_{ij}$ denotes the Kronecker symbol; see, e.g., \cite{krantz1982}.
The norm of the Schwarzian derivative operator of $f$ on the polydisk or the
Euclidean unit ball is given by
\[
\|S_f(z)\|=\sup_{\|\vec{v}\|=1}\|S_f(z)(\vec{v},\vec{v})\|\,,
\]
where the norm is the Bergman norm of the respective domain. Finally, we set
\[
\|S_f\| = \sup_{z\in\Omega}\|S_f(z)\|\,,\qquad \Omega=\mathbb{P}^n
\text{ or }\Omega=\mathbb{B}^n.
\]

In this paper, we use these norms to measure the size of the Schwarzian
operator for convex mappings on the polydisk and for the Roper--Suffridge
extension in the Euclidean unit ball of $\mathbb{C}^n$.

\section{Schwarzian Derivative for Convex Mappings}\label{sec:main}

We begin with the polydisk. Let $f:\mathbb{P}^n\to\mathbb{C}^n$ be a convex
mapping. By the structure theorem for convex mappings on the polydisk, $f$
can be written as
\[
f(z)=f(z_1,\ldots,z_n) = M\bigl(\varphi_1(z_1),\dots,\varphi_n(z_n)\bigr),
\]
where each $\varphi_i$ is a convex holomorphic mapping of the unit disk
$\mathbb{D}$, and $M\in\mathcal{L}(\mathbb{C}^n,\mathbb{C}^n)$ is a linear
map. Since $f=M\circ \varphi$, where $M(z)=M\cdot z$ is linear and
$\varphi(z)=(\varphi_1(z_1),\ldots, \varphi_n(z_n))$, the chain rule
\eqref{chain rule} gives $S_f(z)=S_\varphi(z)$. Without loss of generality,
we may therefore assume that $M$ is the identity. Straightforward calculations
then show that
\[
Df(z)=
\begin{bmatrix}
  \varphi_1'(z_1) & 0 & \cdots & 0 \\
  0 & \ddots & & \vdots \\
  \vdots & & \ddots & 0 \\
  0 & \cdots & 0 & \varphi_n'(z_n)
\end{bmatrix}.
\]
The Jacobian of $f$ is given by
$J_f(z)=\varphi_1'(z_1)\cdots\varphi_n'(z_n)$, which implies that
\[
\nabla\log J_f(z)
=\Bigl(\dfrac{\varphi_1''(z_1)}{\varphi_1'(z_1)},\dots,
\dfrac{\varphi_n''(z_n)}{\varphi_n'(z_n)}\Bigr).
\]
Thus, the Schwarzian derivative tensor applied to the vector
$\vec{v}=(v_1,\ldots,v_n)$ is given by
\[
S_f(z)(\vec{v},\vec{v})
= D^2f(z)(Df(z))^{-1}(\vec{v},\vec{v})
-\frac{2}{n+1}\bigl(\nabla\log J_f(z)\cdot\vec{v}\bigr)\,\vec{v},
\]
or equivalently,
\[
S_f(z)(\vec{v},\vec{v})
=
\begin{pmatrix}
  \displaystyle\frac{\varphi_1''}{\varphi_1'}(z_1)v_1^2
  -\displaystyle\frac{2}{n+1}
  \Bigl(\sum_{j=1}^n\frac{\varphi_j''}{\varphi_j'}(z_j)v_j\Bigr)v_1 \\
  \vdots \\
  \displaystyle\frac{\varphi_n''}{\varphi_n'}(z_n)v_n^2
  -\displaystyle\frac{2}{n+1}
  \Bigl(\sum_{j=1}^n\frac{\varphi_j''}{\varphi_j'}(z_j)v_j\Bigr)v_n
\end{pmatrix}
=
\begin{pmatrix} \delta_1 v_1\\\vdots \\ \delta_n v_n\end{pmatrix},
\]
where
\[
\delta_k(z)=\frac{\varphi_k''}{\varphi_k'}(z_k)v_k
-\frac{2}{n+1}\sum_{j=1}^n\frac{\varphi_j''}{\varphi_j'}(z_j)v_j.
\]

\begin{theorem}
Let $f:\mathbb{P}^n\to\mathbb{C}^n$ be a locally univalent convex mapping.
Then
\[
\|S_f\|\leq \frac{\sqrt{8(n+3)(n-1)}}{n+1}.
\]
The bound is sharp.
\end{theorem}

\begin{proof}
Let $\vec{v}$ be a unit vector in the Bergman metric of the polydisk, so that
\[
\|\vec{v}\|=\left[2\sum_{i=1}^n\frac{|v_i|^2}{(1-|z_i|^2)^2}\right]^{1/2}=1.
\]
Then, since $S_f(z)(\vec{v},\vec{v})=(\delta_k v_k)_k$ and the Bergman norm
of this vector is bounded by $\max_k|\delta_k|\cdot\|\vec{v}\|$, we have
$\|S_f(z)\|\leq m(z)\|\vec{v}\|=m(z)$, where
$m(z)=\max\{|m_k(z)|:k=1,\ldots,n\}$ and
\[
m_k(z) =\left(\frac{-2}{n+1}\frac{\varphi_1''}{\varphi_1'}(z_1)v_1
\right)
+\cdots+\left(\frac{n-1}{n+1}\frac{\varphi_k''}{\varphi_k'}(z_k)v_k
\right)
+\cdots+\left(\frac{-2}{n+1}\frac{\varphi_n''}{\varphi_n'}(z_n)v_n
\right).
\]
Rewriting $m_k$ as
\[
|m_k(z)|=
\left|a_1\cdot\frac{v_1}{1-|z_1|^2}+\cdots+a_n\frac{v_n}{1-|z_n|^2}\right|,
\]
where
\[
a_i=\frac{-2}{n+1}\frac{\varphi_i''}{\varphi_i'}(z_i)(1-|z_i|^2),
\quad i\neq k,
\qquad
a_k=\frac{n-1}{n+1}\frac{\varphi_k''}{\varphi_k'}(z_k)(1-|z_k|^2),
\]
and applying the Cauchy--Schwarz inequality, we obtain
\[
|m_k(z)| \leq
\left(\sum_{i=1}^{n}|a_i|^2\right)^{1/2}
\left(\sum_{i=1}^{n}\frac{|v_i|^2}{(1-|z_i|^2)^2}\right)^{1/2}
=\left(\sum_{i=1}^{n}|a_i|^2\right)^{1/2}\frac{1}{\sqrt{2}}.
\]
Since each mapping $\varphi_i$ is convex, the standard estimate
$|\varphi_i''/\varphi_i'(z_i)|(1-|z_i|^2)\leq 4$ holds, which implies that
\[
\left(\sum_{i=1}^{n}|a_i|^2\right)^{1/2}
\leq
\left(\frac{4}{(n+1)^2}\cdot16(n-1)+16\left(\frac{n-1}{n+1}\right)^2\right)^{1/2}
=\left(\frac{16(n-1)(n+3)}{(n+1)^2}\right)^{1/2}.
\]
This yields
\[
\|S_f\|\leq \frac{2\sqrt{2(n+3)(n-1)}}{n+1}.
\]

Sharpness is achieved for $f=(\varphi_1,\ldots,\varphi_n)$ with
$\varphi_i(z_i)=1/(1-z_i)$ for all $i=1,\ldots,n$. In this case,
\[
|a_i|=\left|\frac{4(1-|z_i|^2)}{(n+1)(1-z_i)}\right|\leq \frac{4(1+|z_i|)}{n+1},
\qquad
a_k=\frac{2(n-1)(1-|z_k|^2)}{(n+1)(1-z_k)},
\]
and equality holds when $z_i\in[0,1)$; moreover, $|a_i|\to 8/(n+1)$ and
$|a_k|\to 4(n-1)/(n+1)$ as $z_i$ and $z_k$ tend to $1$ along the real axis.
For any $z\in\mathbb{P}^n$ one may take
\[
v_i=\frac{\overline{a_i}(1-|z_i|^2)}{\sqrt{2(|a_1|^2+\cdots+|a_n|^2)}},
\quad i=1,\ldots,n,
\]
so that
\[
\lim_{z_i\to 1^-}m_k(z)
=\lim_{z_i\to 1^-}\frac{\sqrt{|a_1|^2+\cdots+|a_n|^2}}{\sqrt{2}}
= \frac{2\sqrt{2(n+3)(n-1)}}{n+1}. \qedhere
\]
\end{proof}

It is worth pausing to check consistency with the one-variable theory. For
$n=2$, the bound gives $2\sqrt{10}/3\approx 2.1081$, and one can verify that
$\|S_f\|<\sqrt{8}\approx 2.8284$ for all $n>1$. The case $n=1$ is more
subtle: returning to the definition of $S^k_{ij}f$ at the beginning of
Section~\ref{sec:prelim}, one sees that for $n=1$ the symmetry condition on
the indices forces $S_f\equiv 0$ identically, for any locally univalent
function. The bound therefore vanishes, which is entirely natural: in one
variable, the tensor-valued Schwarzian is trivial, and the classical
one-variable Schwarzian is a genuinely different object.

\subsection{The Roper--Suffridge Operator}

Working directly with convex mappings in the Euclidean unit ball $\mathbb{B}^n$
is considerably more delicate than in the polydisk. Although a convexity
criterion for mappings on $\mathbb{B}^n$ exists, it involves conditions on the
full Jacobian matrix that do not lend themselves to explicit norm calculations
for $S_f$. The Roper--Suffridge operator sidesteps this difficulty in an
elegant way: by lifting a convex function from the disk to the ball via a
natural geometric construction, it produces a rich and explicit family of
convex mappings in $\mathbb{B}^n$. Within this class, the Schwarzian operator
takes a particularly tractable form, and a precise norm bound becomes
accessible. The estimate we obtain is partial in scope --- it applies only to
this class of mappings --- but it represents the best available result in the
ball setting.


Let $z\in \mathbb{D}$ and let $\varphi$ be a locally univalent function in
$\mathbb{D}$ of the form $\varphi(z)=z+a_2z^2+\cdots$. The Roper--Suffridge
extension operator is defined by
\begin{equation}\label{roper-suff}
\Phi(\varphi)(z_1,z_2)=f(z_1,\ldots,z_n)
=\left(\varphi(z_1),\,\sqrt{\varphi'(z_1)}\,z_2,
\ldots,\sqrt{\varphi'(z_1)}\,z_n\right),
\end{equation}
where the branch of the square root is chosen so that $\sqrt{\varphi'(0)}=1$.
Note that if $\varphi$ is univalent in $\mathbb{D}$, then $\Phi(\varphi)$ is
univalent in $\mathbb{B}^n$. This operator was introduced in \cite{RS} in order
to construct convex mappings in the Euclidean ball in $\mathbb{C}^n$ from
convex functions in the unit disk; that is, if $\varphi$ is a convex function
in $\mathbb{D}$, then $\Phi(\varphi)$ is a convex mapping in $\mathbb{B}^n$.
The Jacobian matrix of $f=\Phi(\varphi)$ is
\[
Df =\begin{pmatrix}
  \varphi' & 0 & \cdots & 0 \\
  \frac{1}{2}z_2\frac{\varphi''}{\sqrt{\varphi'}} & \sqrt{\varphi'} & \cdots & 0 \\
  \vdots & \vdots & \ddots & \vdots \\
  \frac{1}{2}z_n\frac{\varphi''}{\sqrt{\varphi'}} & 0 & \cdots & \sqrt{\varphi'}
\end{pmatrix},
\]
and the Jacobian determinant is $J_f(z)=(\varphi'(z_1))^{(n+1)/2}$. A direct
computation shows that the matrices of the Schwarzian derivative $S^kf(z)$ are
given by
\[
S^kf(z)=\begin{pmatrix}
  \frac{1}{2}z_k S_\varphi(z_1) & 0 & \cdots & 0 \\
  0 & 0 & \cdots & 0 \\
  \vdots & \vdots & \ddots & \vdots \\
  0 & 0 & \cdots & 0
\end{pmatrix},
\qquad k=2,\ldots,n,
\]
and $S^1f(z)$ is the zero matrix. Moreover, for any
$\vec{v}=(v_1,\ldots,v_n)$,
\[
S_f(z)(\vec{v},\vec{v}) =
\begin{pmatrix}
  0 \\
  \frac{1}{2}z_2 S_\varphi(z_1)v_1^2 \\
  \vdots \\
  \frac{1}{2}z_n S_\varphi(z_1)v_1^2
\end{pmatrix}.
\]

\begin{theorem}
Let $\varphi$ be a convex holomorphic function in $\mathbb{D}$ and let
$f = \Phi(\varphi)$ be its Roper--Suffridge extension to $\mathbb{B}^n$.
Then
\[
\|S_f\| \leq \frac{2}{3\sqrt{3(n+1)}}.
\]
\end{theorem}
\begin{proof}
The norm of this operator (computed with respect to the Bergman metric) is
\[
\|S_f(z)(\vec{v},\vec{v})\|^2
=\sum_{i,j=2}^n\frac{n+1}{(1-|z|^2)^2}
\left[(1-|z|^2)\delta_{ij}+\overline{z_i}z_j\right]
\frac{1}{4}z_i\overline{z_j}|S_\varphi(z_1)|^2|v_1|^4.
\]
Since $\|\vec{v}\|=1$, the norm of $S_f(z)$ is attained at
$\vec{v}=(v_1,0,\ldots,0)$, which gives
\[
\frac{|v_1|^2(n+1)(1-|z|^2+|z_1|^2)}{(1-|z|^2)^2}=1.
\]
A straightforward calculation yields
\[
\sum_{i,j=2}^n\left[(1-|z|^2)\delta_{ij}+\overline{z_i}z_j\right]z_i\overline{z_j}
=(1-|z_1|^2)\sum_{i=2}^n|z_i|^2,
\]
and hence
\[
\|S_f(z)(\vec{v},\vec{v})\|^2
=\frac{|S_\varphi(z_1)|^2(1-|z_1|^2)^4}{4(n+1)}
\cdot\frac{(1-|z|^2)^2\displaystyle\sum_{i=2}^n|z_i|^2}
{(1-|z_1|^2)^3\!\left(1-\displaystyle\sum_{i=2}^n|z_i|^2\right)^{\!2}}.
\]
A direct optimization shows that
\[
\sup\left\{\frac{(1-x-y)^2y}{(1-x)^3(1-y)^2}
:x\geq 0,\;y\geq 0,\;x+y<1\right\}=\frac{4}{27},
\]
where the supremum is not attained but is approached in the limit as
$(x,y)\to(1,0)$.
Setting $x=|z_1|^2$ and $y=\sum_{i=2}^n|z_i|^2$, we obtain
\[
\sup_{z\in\mathbb{B}^n}
\frac{(1-|z|^2)^2\displaystyle\sum_{i=2}^n|z_i|^2}
{(1-|z_1|^2)^3\!\left(1-\displaystyle\sum_{i=2}^n|z_i|^2\right)^{\!2}}
=\frac{4}{27},
\]
with the supremum attained in the limit as $z_1\to 1$. Therefore,
\[
\|S_f\|=\frac{1}{3\sqrt{3(n+1)}}
\sup_{z_1\in\mathbb{D}}|S_\varphi(z_1)|(1-|z_1|^2)^2.
\]
Since $\varphi$ is a convex mapping, we conclude that
\[
\|S_f\| \leq \frac{2}{3\sqrt{3(n+1)}}.
\]

\end{proof}

A comprehensive treatment of the Roper--Suffridge operator, including its
role in the construction of starlike and convex mappings in higher dimensions,
can be found in \cite[Ch.~11]{graham_kohr2003}.



% --- Bibliografía ---
\bibliographystyle{amsplain}
\bibliography{schwarziana}

\end{document}
